\documentclass[12pt]{article}
\usepackage[T1]{fontenc}
\usepackage{graphicx,amsmath,amssymb}
\usepackage[margin=25mm]{geometry}
\usepackage{setspace}
\usepackage[super,sort&compress]{natbib}
\usepackage{lineno}
\usepackage{booktabs,array}
\usepackage[hidelinks]{hyperref}
\doublespacing
\raggedright
\setlength{\parindent}{0pt}
\setlength{\parskip}{5pt}
\begin{document}
\title{Bayesian inference of target-associated migration changes in synthetic killer-cell trajectories}
\author{}
\date{}
\maketitle
\begin{center}\small Computational research manuscript draft. All results are synthetic.\end{center}
\begin{abstract}
Target-associated changes in immune-cell movement are difficult to distinguish from persistent random motion, mechanical retention and imaging noise. Here we develop a synthetic benchmark combining an accelerated agent-based simulator, hierarchical Bayesian state-space inference and neural simulation-based inference. Both approaches recover strong proximity-associated reductions in active movement and increases in short-timescale jitter, whereas persistence is less identifiable. On independent simulations, neural 90\% intervals achieve @@NEURAL_COVERAGE_RANGE@@ coverage. Expanded tests show @@NEURAL_NULL_COUNT@@ declarations among 256 unchanged-motility movies. A model trained at low imaging noise incorrectly declares an effect in all 128 high-noise controls; broader noise training reduces this to @@NOISE_MARG_COUNT@@, and conditioning on known noise to @@NOISE_COND_COUNT@@. Mechanical and temporal misspecification can nevertheless bias parameters without conspicuous predictive failure. These results establish recoverability under explicit assumptions and provide a reproducible framework for evaluating target-associated migration hypotheses before experimental interpretation.
\end{abstract}
\clearpage
\linenumbers

Immune-cell migration determines which potential targets are encountered, how long cells remain in their vicinity and whether successive interactions can occur. A trajectory that slows near a target and displays frequent short displacements can suggest retention or local searching. Yet this description combines distinct properties of motion. Smaller persistent displacements, faster loss of directional memory, stronger short-timescale fluctuations and adhesive constraints can all contribute to a visually irregular path. Microscopy adds another source of apparent movement: localization errors become differenced when velocity is estimated from recorded positions. Establishing which of these explanations is supported by a trajectory is therefore an inference problem, rather than a property that can be read directly from a track.

Work on heterogeneous random walks has shown that population-averaged migration statistics can conceal variation among cells and within individual trajectories. Superstatistical analysis inferred time-varying persistence and activity, including gradual and abrupt changes within heterogeneous trajectories\citep{metzner2015}. A subsequent Bayesian framework extended parameter-evolution modeling and comparison between dynamic hypotheses\citep{mark2018}. These approaches motivate an analysis at the level of local motion, rather than a single mean speed or mean-squared displacement. For killer--target interactions, however, identifying a temporal change is only part of the question. A target-distance explanation must also be distinguished from unrelated temporal variation and from the physical consequences of contact. Experimental observations of natural killer-cell scanning and prolonged target attachment provide biological motivation for testing these alternatives, without determining which effective stochastic parameters should change\citep{olofsson2014}. Work on target-directed interactions also illustrates why an encounter alone is insufficient evidence of purposeful pursuit: directional behavior must be compared with appropriate target-based reference distributions\citep{metzner2021}.

Bayesian state-space models provide a natural route from noisy positions to uncertainty in latent motion. For a persistent Gaussian velocity process, latent states can be integrated out sequentially using a Kalman filter\citep{kalman1960}. A hierarchical model can then accommodate differences among cells while estimating a shared response to target proximity. Simulation-based inference offers a complementary route when a simulator is easier to construct than a tractable likelihood\citep{cranmer2020,papamakarios2016}. Neural conditional-density estimators can make repeated inference inexpensive after training, but their uncertainty and response to model mismatch require explicit validation\citep{lueckmann2021}. Agreement between fitted and simulated summaries is informative only for the features, parameter ranges and observation conditions represented in the evaluation.

Here we construct a controlled computational study of target-associated motility inference. The aim is to determine which effective changes can be recovered, how two inference approaches compare and which departures from their assumptions produce misleading conclusions. An accelerated agent-based model generates killer trajectories around known target positions with separately controlled active movement amplitude, persistence and translational jitter. A hierarchical state-space model and a neural posterior estimator analyze the same observation schema, with rejection approximate Bayesian computation providing a simple likelihood-free comparator. Independent simulations test interval coverage, parameter recovery and prespecified directional decisions. Follow-up analyses expand replication, vary baseline parameters and neural training seeds, and evaluate broader training over imaging noise. An independently implemented supplied agent-based model and explicit mechanical and temporal stress cases assess transfer beyond the principal training process. The contribution is a reproducible account of recoverability and failure modes; the study contains no experimental trajectories and makes no claim that real killer cells regulate migration near targets.

\section*{Results}
\subsection*{A controlled model separates three motility responses}
We represented each killer by a two-dimensional latent velocity and position. Active velocity followed an Ornstein--Uhlenbeck process with local stationary per-axis standard deviation $S$ and persistence time $\tau$; independent Brownian position increments had amplitude $J$. These parameters describe different components of movement. In particular, $S$ is not the mean observed scalar speed, and an increase in $J$ can increase frame-to-frame displacement even when persistent movement decreases. Target proximity was encoded by a continuous gate that equals one at the assumed contact distance and decays exponentially with positive surface gap (Fig.~1a,b). The logarithms of $S$, $\tau$ and $J$ varied linearly with this gate. A cell-specific random intercept affected $S$, enabling within-cell distance variation to contribute information beyond differences among cells.

The benchmark used 12 killers, nine stationary targets and 201 observations at 0.25-min intervals in a periodic 180-$\mu$m square. Target positions were known without measurement error; independent Gaussian localization noise affected killer positions. These are synthetic design choices, not measurements of a biological preparation. The baseline geometry allowed target zones to be crossed because exclusion forces were absent. It therefore tested inference of proximity-dependent movement. Separate mechanical simulations added soft repulsion and stochastic adhesion to test the consequences of omitted contact physics.

The strong combined alternative reduced active amplitude and persistence to approximately 0.45 times their far-field values and increased jitter to approximately 2.23 times its far-field value at contact. Isolated alternatives changed one component at a time, while a weaker combined alternative used smaller log effects. The unchanged-motility condition set all three proximity coefficients to zero. Null trajectories still contained ordinary stochastic motion, cell heterogeneity and repeated target encounters. No movies were selected on the basis of whether a cell reached, remained near or moved away from a target.

The two inference routes estimated the same six global parameters (Fig.~1c). The state-space model also sampled 12 cell intercepts while integrating latent position and velocity analytically. Neural inference conditioned on 82 observed-track summaries and marginalized cell heterogeneity through simulation. Both methods used the same bounded parameter prior. Evaluating posterior response functions along observed distance histories yielded time-dependent parameter paths (Fig.~1d--f). These paths share a common distance-response model; they are not independent estimates of three new parameters at every frame.

\subsection*{Accelerated simulation supports independent validation}
The reduced simulator used compiled numeric kernels and the exact joint Gaussian transition of velocity and integrated position when local coefficients were held fixed within a timestep. Retaining the velocity--position noise covariance avoided treating integrated motion as independent of the updated velocity. A reference Python engine received identical random draws and produced identical recorded coordinates in the paired timing test. After warm-up, the compiled engine required 0.327\,ms per movie, compared with 61.84\,ms for its equivalent Python engine, a 189.2-fold speedup. These timings describe the stated small-population workload; they do not establish the same speedup for dense populations, different hardware or the complete supplied agent-based model.

A 16,000-movie prior bank, including trajectory summaries, was generated in 43.3\,s, and the original neural model trained in 9.7\,s. This computational cost made it practical to reserve complete, independently generated movies for evaluation. The expanded study used 1,024 new prior-predictive movies for neural calibration and 96 new movies for state-space calibration. Additional fixed-null, varied-baseline-null and alternative simulations evaluated decisions independently of the calibration ensemble. The follow-up plan was written before those simulations were run, but was informed by the initial benchmark and was not a preregistration (Table~1).

Numerical validation addressed the implementation as well as the statistical output. Tests compared the integrated-velocity covariance with matrix-exponential and numerical-integration calculations, compared the Kalman likelihood with a full Gaussian calculation, checked missing-observation handling and verified that replacing stored truth labels did not affect inference. Source-ABM regression tests preserved the original seeded behavior when the new proximity effects were zero. All 41 core tests passed. Every newly reported state-space ensemble retained its convergence attempts and unresolved fits rather than removing difficult cases from the record (Supplementary Fig.~S7).

\subsection*{Calibration accompanies useful parameter information}
For the frozen original neural estimator, nominal 90\% marginal intervals achieved @@NEURAL_COVERAGE_RANGE@@ empirical coverage across the six parameters in the 1,024-movie follow-up ensemble (Fig.~2a; Supplementary Fig.~S1). The intervals were narrower than the corresponding central 90\% prior intervals, indicating that coverage was accompanied by information gain (Fig.~2b). Posterior-median error was smallest for the baseline jitter scale and larger for the proximity effect on persistence. The root-mean-squared errors for the three proximity effects were @@NEURAL_EFFECT_RMSE@@ in log-parameter units. The persistence-effect interval retained @@TAU_RELATIVE_WIDTH@@ of the prior interval width, compared with @@S_RELATIVE_WIDTH@@ for active amplitude and @@J_RELATIVE_WIDTH@@ for jitter.

The new state-space calibration ensemble produced @@SSM_CAL_PASSED@@ numerically accepted fits out of 96 attempts. Marginal 90\% coverage across parameters ranged from @@SSM_COVERAGE_RANGE@@, with the individual binomial uncertainty intervals reported in Supplementary Table~S2. Its proximity-effect root-mean-squared errors were @@SSM_EFFECT_RMSE@@. The likelihood-based and neural approaches therefore provided useful information under their respective assumptions, while differing in the amount of information retained about persistence (Fig.~2c--e). State-space coverage is evaluated against a simulator whose distance gate is computed from latent positions, whereas the fitted gate is computed from noisy observations. Its likelihood is exact for the conditional model, but this difference prevents describing the entire simulator-to-likelihood comparison as analytically exact.

Three additional neural networks trained using independently seeded initializations and training/validation splits of the same 16,000-movie bank. Their coverage ranges were @@SEED_COVERAGE_RANGES@@ (Supplementary Fig.~S2). Reusing a bank holds its particular simulation realization fixed; this experiment measures training variability, not variability across independently generated training banks. The original network was frozen before the expanded evaluation, and no evaluation result was used to choose a winning model. All network results are reported separately, because the same test movies evaluated by several networks are not additional independent datasets.

Rejection approximate Bayesian computation retained the nearest 256 parameter draws from the original 16,000-movie bank in standardized summary space. In the initial independent 256-movie test, its 90\% intervals covered approximately 92.6--95.7\% of generating values, but were generally broader and had larger errors than the neural estimator. This conservative comparator illustrates why coverage alone is insufficient: a diffuse posterior can cover the truth while distinguishing relatively few alternatives. Its nonzero tolerance and correlated summary features also prevent treating it as an exact reference posterior.

\subsection*{Persistence is less identifiable than slowdown or jitter}
We declared an individual effect only when its marginal posterior probability exceeded 0.95 of crossing a prespecified 1.2-fold ratio threshold in the hypothesized direction. The three directions were smaller active amplitude, shorter persistence and larger jitter. An any-effect declaration was the logical union of those three separate decisions. This rule is not a Bayes factor or a fixed 5\% frequentist test, and absence of a declaration is inconclusive rather than evidence of equivalence.

The frozen neural estimator declared an effect in @@NEURAL_NULL_COUNT@@/256 fixed-baseline null movies and @@NEURAL_VARIED_NULL_COUNT@@/256 movies with null proximity effects and baseline parameters drawn across the prior. The corresponding 95\% Wilson intervals were @@NEURAL_NULL_WILSON@@ and @@NEURAL_VARIED_NULL_WILSON@@ (Fig.~3a). State-space inference declared effects in @@SSM_NULL_COUNT@@/@@SSM_NULL_N@@ accepted fixed-null fits and @@SSM_VARIED_NULL_COUNT@@/@@SSM_VARIED_NULL_N@@ accepted varied-baseline-null fits, with uncertainty intervals @@SSM_NULL_WILSON@@ and @@SSM_VARIED_NULL_WILSON@@. These results describe the tested nuisance settings and parameter distributions; they do not bound false declarations under arbitrary imaging or mechanical conditions.

The varied-baseline null also exposed a conditional calibration limitation. Its baseline-persistence parameter was covered by the state-space 90\% interval in only @@SSM_VARIED_TAU_COUNT@@/48 movies (@@SSM_VARIED_TAU_COVERAGE@@; 95\% Wilson interval @@SSM_VARIED_TAU_WILSON@@). Thus, satisfactory prior-predictive coverage averaged across generating effects did not ensure the same coverage in this null subset. The three varied-null effect declarations all concerned persistence.

In the strong combined alternative, neural inference recovered slowdown in @@NEURAL_COMBINED_S_COUNT@@/128 movies and increased jitter in @@NEURAL_COMBINED_J_COUNT@@/128, but shorter persistence in only @@NEURAL_COMBINED_T_COUNT@@/128 (Fig.~3b,c). State-space inference detected the three effects in @@SSM_COMBINED_S_COUNT@@, @@SSM_COMBINED_T_COUNT@@ and @@SSM_COMBINED_J_COUNT@@ of @@SSM_COMBINED_N@@ accepted combined-condition fits, respectively. The independently generated ensembles support comparison within the same design, while a paired evaluation of the neural models on the exact state-space movies is provided in the source data and Supplementary Information. In the persistence-only alternative, the frozen neural model detected the change in @@NEURAL_SOLO_T_COUNT@@/128 movies. This larger follow-up estimate was lower than the initial 17/32 estimate, illustrating the sampling uncertainty of the exploratory benchmark.

The reduced sensitivity to persistence has a direct model-based explanation, although this does not constitute an identifiability proof. For locally constant coefficients and stationary velocity, the per-axis displacement variance over interval $h$ is
\begin{equation}
\operatorname{Var}(\Delta y)=2S^2\tau\{h-\tau(1-e^{-h/\tau})\}+J^2h+2\sigma_{\rm obs}^2.
\end{equation}
At short intervals, the active contribution is $S^2h^2-S^2h^3/(3\tau)+O(h^4)$, whereas translational jitter contributes at order $h$ and localization noise contributes a constant. Persistence must therefore be inferred largely from temporal structure that can become less visible when jitter increases. At the chosen baseline and recording interval, the Brownian component accounts for approximately 16\% of the locally stationary displacement variance far from targets and 79\% at contact under the combined alternative (Supplementary Fig.~S8). These are analytical illustrations at zero cell offset, not observed biological fractions. A positive any-effect decision consequently cannot be interpreted as successful recovery of every programmed change.

\subsection*{Broader noise training repairs a major failure mode}
Imaging noise was the strongest failure in the initial neural stress test. The original estimator was trained with localization standard deviation 0.15\,$\mu$m. When the unchanged-motility simulator instead used 0.60\,$\mu$m, that estimator declared an effect in all 128 new null movies (Fig.~4a). This was a deliberately unguarded extrapolation test. The standard application wrapper rejects a declared mismatch in acquisition settings; it cannot protect against noise that is unknown or incorrectly specified.

To investigate whether the failure was remediable, we generated a new 24,000-movie bank in which localization standard deviation varied log-uniformly from 0.05 to 0.70\,$\mu$m. A noise-conditioned estimator received the logarithm of the known standard deviation alongside the 82 track summaries. A matched-bank ablation used the same expanded simulations without this additional scalar, thereby marginalizing noise rather than conditioning on its known value. The first specified seed in each family was designated the primary model before evaluation; two further seeds assessed training sensitivity.

The primary noise-conditioned model declared an effect in @@NOISE_COND_COUNT@@/128 high-noise null movies and @@NOISE_COND_LOW_COUNT@@/128 low-noise null movies. Across 1,024 independent prior-predictive movies spanning the expanded noise range, its nominal 90\% coverage was @@NOISE_COND_COVERAGE@@ (Fig.~4b). Strong combined changes remained detectable in @@NOISE_COND_COMBINED_COUNT@@/128 high-noise movies, although individual persistence sensitivity remained lower than that for the other effects (Fig.~4c). The high-noise null bias in the active-amplitude response was also reduced (Fig.~4d).

Importantly, broader noise training without an explicit noise input already removed much of the failure: the primary marginalized model declared effects in @@NOISE_MARG_COUNT@@/128 high-noise controls and detected any effect in @@NOISE_MARG_COMBINED_COUNT@@/128 high-noise combined movies. Its coverage range was @@NOISE_MARG_COVERAGE@@. Thus, the intervention supports matching the training distribution to the observation process; it does not establish that the extra known-noise scalar alone caused the improvement. The three conditioned seeds produced @@NOISE_COND_SEED_COUNTS@@ high-noise null declarations, and the corresponding marginalized counts were @@NOISE_MARG_SEED_COUNTS@@ (Supplementary Fig.~S3). This experiment addresses localization noise within a specified range, not noise processes or ranges absent from training, target-position error, blur, tracking swaps or changes in cell density.

\subsection*{Transfer and mechanical tests limit attribution}
The supplied agent-based model provided an independent generator with a fine-step Euler polarity process, zero initial polarity and reciprocal mechanical interactions. Optional proximity coefficients and cell-specific speed factors were added while preserving the original zero-effect dynamics. The bridge disabled killing and exhaustion to focus on migration. Its movement process is not identical to the reduced velocity model when the displacement coupling changes with distance, making this a transfer stress test rather than a second matched calibration experiment.

For the combined alternative generated by that model, neural inference declared some effect in 16/16 movies. The per-effect detections were 14/16 for slowdown, 0/16 for shorter persistence and 16/16 for increased jitter (Fig.~5a). A four-movie state-space pilot detected the same effects in 4/4, 1/4 and 4/4 movies. Both methods detected altered motion more readily than they recovered all its effective components. These small samples and structural differences do not support accurate parameter transfer across arbitrary agent-based models.

Mechanical-only simulations had zero programmed proximity effects on the intrinsic stochastic coefficients, but adhesion and repulsion genuinely changed movement near targets. Positive intrinsic-effect declarations in this setting indicate attribution to an incomplete model, rather than detection of movement that does not exist. Neural inference declared effects in 2/32 such movies, while the state-space pilot declared an effect in 1/8 (Fig.~5b). More revealingly, the state-space 90\% interval for the active-amplitude proximity coefficient covered its zero generating value in only 2/8 movies. With spontaneous temporal fluctuations independent of target distance, it covered zero in only 3/8 despite no positive directional declarations (Fig.~5c). Thus, a low count of practical-threshold decisions can coexist with biased and poorly calibrated parameter estimates.

The selected summary-based posterior predictive checks did not reliably reveal this problem. Only 2.4\% and 3.7\% of summaries fell outside pointwise 90\% predictive intervals for the first mechanical example under neural and state-space inference, respectively (Supplementary Fig.~S6). The summaries are correlated and were not an independent test set; this descriptive agreement cannot validate a mechanism. Mechanical integration was also more sensitive to timestep than the baseline process. One internal step per frame differed from the 16-step reference by approximately 4.9 combined Monte Carlo standard errors for turning and 5.6 for overlap, while four steps retained a turning difference of approximately 2.1 standard errors (Supplementary Fig.~S5). Mechanical conclusions are therefore exploratory and require finer-step and broader mechanism validation before quantitative use.

\section*{Discussion}
This study shows that strong target-associated reductions in persistent movement amplitude and increases in translational jitter can be recovered from synthetic trajectories using both a hierarchical likelihood model and neural simulation-based inference. It also identifies a boundary of that conclusion: successful detection of altered movement does not establish accurate recovery of every motility parameter or identify the physical mechanism responsible. Persistence was less informative than amplitude or jitter, and omitted mechanics and temporal variation could distort coefficient estimates even when a prespecified decision remained negative. The most severe observation-noise failure was substantially reduced by broader simulator training, demonstrating a practical route to improving inference within a defined observation model.

The framework builds on the local, heterogeneous view of migration developed in superstatistical and Bayesian trajectory analyses\citep{metzner2015,mark2018}. Its principal change in scientific emphasis is the explicit association with target geometry and the controlled evaluation of competing explanations. A single-cell trajectory supplies a sequence of motion increments, but its interpretation also depends on which environments the cell visits. A shared distance-response function combines information across encounters while retaining a cell-specific baseline. This is a useful restricted model when the question concerns proximity effects. It does not replace models of arbitrary temporal parameter drift, latent behavioral states or history-dependent responses. In particular, the present inference is retrospective: the likelihood updates latent states sequentially, while global parameters are inferred using complete movies.

The state-space and neural routes have complementary advantages. The former provides an explicit observation model and a traceable likelihood approximation, and retained more useful persistence information in the combined scenario. The latter amortizes the cost of inference across simulations and accommodates an expanded observation process without deriving a new trajectory likelihood. The results do not establish a universal ranking. Neural performance depends on the chosen summaries and density family; richer temporal features or learned sequence representations could alter the gap. State-space performance depends on whether the assumed conditional dynamics are adequate. Providing the correct localization variance to the likelihood also gives it information unavailable to a neural estimator trained at a different fixed noise level, so the original noisy stress test should not be interpreted as an equal-information superiority comparison.

The noise intervention emphasizes the distinction between training coverage and scientific robustness. Both the conditioned and marginalized estimators were exposed to a broader observation distribution. Their improvement shows that the original failure was not an unavoidable consequence of using simulation-based inference, but also that nominal uncertainty from an accurately trained narrow model need not generalize outside its design. Explicitly providing a known nuisance value can be useful, yet it requires that the value itself be calibrated. An uncertain or biased localization estimate would need its own model or sensitivity analysis. Neither an input-range warning nor posterior predictive agreement can guarantee detection of an unknown mismatch. Explicit discrepancy modeling provides a further methodological direction when simulator misspecification cannot be removed\citep{ward2022}.

Biological interpretation requires additional restraint. Physical attachment is itself a biologically meaningful behavior, but it is not equivalent to a change in the coefficients of an unconstrained random walk. Experimental observations of target scanning and prolonged interactions motivate the mechanical alternatives\citep{olofsson2014}; they do not imply that a fitted decrease in $S$ or $\tau$ identifies a signaling pathway. To distinguish active regulation from adhesion, a future generative model should represent both mechanisms and evaluate whether their predictions are distinguishable at the actual imaging cadence. Even then, movement-based association would require suitable perturbations or experimental contrasts before a causal sensing interpretation became justified.

Several features limit the current study's domain. The benchmark is two-dimensional and periodic, uses fixed-size circular proximity zones and assumes noiseless target positions. It omits nonperiodic boundaries, deforming masks, motion blur, missing identities, track fragmentation, target death and reciprocal mechanical coupling in the reduced model. The hierarchy varies only active amplitude and fixes its population standard deviation. Localization variance, interaction length and contact distance are known in each main fit, rather than jointly inferred. The asymmetric effect prior retains zero and both directions but encodes a synthetic design choice that requires sensitivity analysis for an experimental application. Broadening the parameter distribution in calibration does not substitute for testing these alternative structures.

The study also has inferential limits. Marginal prior-predictive coverage and rank histograms do not establish joint or conditionally calibrated uncertainty\citep{talts2018,sailynoja2022}. The repeated neural fits share simulation banks and evaluation movies; their results measure variation over training rather than independent scientific replication. The main large follow-up ensembles were planned after inspecting exploratory results, and the source-ABM and mechanical samples remain small. A finite number of null movies constrains a declaration rate only for the evaluated design. Numerical convergence checks diagnose sampling quality, not model validity. These distinctions are central to a simulation study intended to guide experimental interpretation\citep{morris2019}.

The practical outcome is a tested starting point for asking whether migration depends on target proximity. For new imaging data, the first steps are to match cadence, geometry and localization error; define distance and contact from the relevant segmentation; retain non-encounters; and check that individual cells provide both near- and far-target information. Mechanism-rich simulations and independent experimental movies can then assess whether inferred changes survive plausible alternatives. The present results support pursuing that analysis for slowdown and jitter, identify persistence as a demanding component, and provide the code, source data and failure cases needed to decide where further model development is warranted.

\section*{Methods}
\subsection*{Study design and analysis chronology}
The study used simulated trajectories exclusively. Its exploratory benchmark established model implementations and identified weak persistence recovery and observation-noise sensitivity. The follow-up computational plans were saved before executing the expanded ensembles, with separate random-seed ranges and predefined numerical criteria. They were informed by the exploratory results and were not preregistered. No experimental specimens, participants, animals or clinical data were used. The independent sampling unit was a complete simulated movie; frames and cells within a movie were not treated as independent evaluation replicates. Table~1 and Supplementary Table~S1 give ensemble sizes and roles; parameter distributions are specified below.

\subsection*{Geometry and proximity response}
Target positions were initialized on a $3\times3$ lattice with independent uniform perturbations of up to 9\,$\mu$m per coordinate and periodic wrapping. Killer positions were sampled uniformly in the 180-$\mu$m square, rejecting initial positions with a target-centroid distance below 12\,$\mu$m. No subsequent encounter selection was applied. The minimum-image distance from killer $i$ to its nearest target at time $t$ was $d_i(t)$. With assumed contact distance $R=12\,\mu$m and interaction length $\ell=10\,\mu$m, the gate was
\begin{equation}
 g_i(t)=\exp\{-\max[d_i(t)-R,0]/\ell\}.
\end{equation}
The nearest target could change over time. Ties selected the first target deterministically. The gate saturated at contact and did not include a separate contact-only jump. True positions defined the simulator gate; observed killer positions and known target positions defined the inference gate.

\subsection*{Stochastic movement and observations}
Using an integrated-velocity construction related to continuous-time correlated random walks\citep{johnson2008}, the baseline model for each Cartesian coordinate was
\begin{align}
 dv_i&=-v_i\,dt/\tau_i+\sqrt{2S_i^2/\tau_i}\,dW_{v,i},\\
 dx_i&=v_i\,dt+J_i\,dW_{x,i}+F_i\,dt,
\end{align}
with independent Brownian drivers across components and cells. Local coefficients obeyed
\begin{align}
 \log S_i&=\mu_S+u_i+\delta_Sg_i,&u_i&\sim N(0,0.25^2),\\
 \log\tau_i&=\mu_\tau+\delta_\tau g_i,&
 \log J_i&=\mu_J+\delta_Jg_i.
\end{align}
The baseline had $F_i=0$. Initial velocity was drawn from its local stationary Gaussian distribution. Killer positions were retained unwrapped for displacement calculations, and periodic distances were used only for interactions. Observations were $y_{it}=x_{it}+\epsilon_{it}$ with independent $N(0,\sigma_{\rm obs}^2I_2)$ errors. The baseline used $\sigma_{\rm obs}=0.15\,\mu$m and a 0.25-min recording interval for 50\,min. Targets were observed without noise.

\subsection*{Exact local transition and computation}
Coefficients were held fixed at the start of each internal integration step. For state $(v,x)$ and interval $h$, define $\phi=e^{-h/\tau}$ and $H=\tau(1-\phi)$. The transition matrix was
\begin{equation}
 A=\begin{pmatrix}\phi&0\\H&1\end{pmatrix},
\end{equation}
and the process covariance entries were
\begin{align}
 Q_{vv}&=S^2(1-\phi^2),\\
 Q_{vx}&=S^2\tau(1-\phi)^2,\\
 Q_{xx}&=S^2\tau^2[2h/\tau-4(1-\phi)+(1-\phi^2)]+J^2h.
\end{align}
A Cholesky factor generated joint velocity and displacement innovations. Stable exponential-difference calculations and a small-step series avoided cancellation. Exactness holds for frozen coefficients; spatial modulation and explicit mechanical forces still require timestep assessment. Numba compiled the loop over frames, killers and targets, with random arrays generated before the kernel. The reference Python engine used the same draws. Benchmarks excluded initial compilation and reported median warm timings. The simplified engine omitted killing, exhaustion, event tables and killer--killer or target--target forces.

\subsection*{Priors and generating scenarios}
The six global parameters were independent uniform draws in their stated coordinates: $\mu_S\in[\log2,\log5]$, $\mu_\tau\in[\log0.8,\log3]$, $\mu_J\in[\log0.4,\log1.2]$, $\delta_S,\delta_\tau\in[-1.2,0.5]$ and $\delta_J\in[-0.5,1.2]$. Thus, baseline scales had log-uniform priors. Fixed scenarios used far-field values $(S,\tau,J)=(3.2,1.6,0.7)$ in the adopted units. The strong combined effects were $(-0.8,-0.8,0.8)$; isolated alternatives changed one corresponding coefficient; the weak alternative used $(-0.3,-0.3,0.3)$. Fixed nulls used zero effects. Varied-baseline nulls drew the three baseline log parameters from the same prior while fixing effects to zero. Cell intercepts were independently regenerated for each movie.

\subsection*{Hierarchical state-space likelihood}
The state-space model conditioned on gates computed from observed positions and used the joint transition above for each axis. Latent velocity and position were marginalized using a Kalman filter\citep{kalman1960}. Initial position was conditioned on its first observation under a flat position prior, leaving uncertainty $\sigma_{\rm obs}^2$; initial velocity had local stationary variance. Missing measurements could skip the observation update, with interpolated positions used only to construct covariates, although the main benchmark had complete tracks. The conditional gate approximation did not integrate out localization error in the distance covariate. The fitted baseline also omitted mechanical forces, spontaneous log-speed variation and within-frame gate changes.

The posterior sampled six global parameters and 12 cell random intercepts. A bounded L-BFGS-B optimization initialized a finite-difference Hessian used to construct proposal scales. Each iteration mixed a Hessian-preconditioned random-walk proposal (probability 0.7) with a multivariate Student-$t$ independence proposal (probability 0.3; seven degrees of freedom; scale factor 1.3). Metropolis--Hastings acceptance included the independence-proposal density ratio. Random-walk covariance and scale adapted during warm-up only. An additional exact conditional update refreshed $\mu_S$ while holding each $\mu_S+u_i$ fixed, improving hierarchical mixing without changing the likelihood. Four independently seeded chains were run. Acceptance rates and optimizer diagnostics were retained.

All global and cell parameters had to satisfy rank-normalized $\widehat R<1.01$ and bulk and tail effective sample sizes of at least 400\citep{vehtari2021}. Follow-up fits started with 2,000 retained iterations per chain and 1,500 warm-up iterations; failing fits were rerun at 4,000, 8,000, 16,000 and at most 32,000 retained iterations, with the warm-up schedule recorded in the analysis plan. Earlier attempts were preserved. Primary metrics included only accepted fits, accompanied by attempted and failed counts and worst/best bounds that treated unresolved outcomes explicitly. This numerical rule did not select on effect direction or recovery accuracy.

\subsection*{Trajectory summaries and neural posterior}
The summary vector contained 82 observed-only features: multiscale displacement statistics, increment correlations, turning, radial motion, distance-stratum exposure, dwell measures and within-cell near/far contrasts. Summary definitions and exact ordering are included in the source snapshot. Neither summaries nor inference read stored latent states, effect labels or generating cell offsets. The neural posterior used three fully covariant Gaussian components over standardized logit-transformed bounded parameters. Three hidden layers of width 128 used SiLU nonlinearities. Cholesky diagonal terms were positive by construction. Density evaluation on the original parameter scale included the transformation Jacobian.

The original 16,000-movie bank supplied a movie-level split with 15\% reserved for internal validation. Training used AdamW with learning rate 0.0015 and weight decay $10^{-4}$, batch size 128, gradient-norm clipping at 5, at most 300 epochs and early-stopping patience 40. Standardized contexts were clipped to $[-12,12]$ inside the network, and Cholesky diagonals used softplus plus 0.025 for numerical stability. Input and parameter transformations were fitted on training rows only. Negative log density was minimized with early stopping on the held-out split, and the best validation checkpoint was restored. The three follow-up fixed-noise seeds reused this bank with independent initialization and split seeds. Scientific evaluation used new simulated movies, not training or early-stopping validation rows. Posterior probabilities and intervals used @@POSTERIOR_DRAWS@@ draws per movie in the expanded neural study. This was a custom PyTorch conditional mixture implementation, not the separately distributed \texttt{sbi} package. It estimates the posterior given the chosen summaries, which can lose information relative to full trajectories\citep{papamakarios2016,lueckmann2021}.

\subsection*{Noise-aware training and matched ablation}
A separate bank of 24,000 movies drew $\sigma_{\rm obs}$ log-uniformly on $[0.05,0.70]\,\mu$m, independently of global parameters and cell effects. All other baseline nuisance settings were fixed. Conditioned networks appended $\log\sigma_{\rm obs}$ to the 82 summaries. Marginalized networks used the same bank and 82 summaries without that scalar. Three seeds per family were trained. Primary seeds were designated before evaluation; remaining fits were sensitivity analyses. A new 1,024-movie ensemble evaluated calibration across noise, and 128 movies per fixed-noise null or combined condition evaluated decisions at 0.15 and 0.60\,$\mu$m. The original fixed-noise model evaluated those same movies as an intentionally mismatched comparator. The intervention assumes known noise for the conditioned model; it does not infer uncertainty in the supplied noise value.

\subsection*{Approximate Bayesian computation}
The rejection comparator standardized features by their prior-bank standard deviations, omitted numerically constant dimensions and ranked bank rows by root-mean-squared standardized Euclidean distance to the observation. It retained the nearest 256 of 16,000 rows without regression adjustment. Accepted rows were actual prior draws. The final accepted distance and acceptance fraction were recorded. Correlated summaries, finite tolerance and finite-bank reuse are limitations of this comparator. Its reported calibration came from the initial independent 256-movie ensemble; it was not included in every expanded comparison.

\subsection*{Stress tests and independent source model}
Spontaneous temporal variation added an independent Ornstein--Uhlenbeck term to log active amplitude, with standard deviation 0.55 and persistence 5\,min. Mechanical-only simulations added Hookean overlap repulsion with coefficient 5\,min$^{-1}$, bond formation rate 2\,min$^{-1}$ within the 12-$\mu$m contact distance, bond dissociation rate 0.5\,min$^{-1}$ and a radial adhesive drift of 5\,$\mu$m\,min$^{-1}$ while bonded. Bonds also broke beyond 1.5 times contact distance. Four internal steps per recorded frame were used in the original mechanical benchmark. Other exploratory stress cases used moving targets or finer movement integration; their complete results remain in the source data.

The supplied MantiShrimp implementation retained its polarity-based Euler dynamics and reciprocal forces. Its optional local effects scaled displacement coupling, polarity decay/noise and translational noise consistently with the intended local response. Zero-effect regression tests checked seeded final states, event sequences and adaptive-step counts. The independent bridge used ten internal steps per recorded frame, zero initial polarity and no killing/exhaustion. Sixteen movies per source-model condition were evaluated by neural and rejection methods, and four by state-space inference. Parameter interpretation across source and reduced models is approximate when coupling varies with distance.

\subsection*{Performance measures and statistical analysis}
Posterior-median error, root-mean-squared error, interval width, marginal coverage at 50\%, 80\%, 90\% and 95\%, and posterior rank fractions were computed against generating parameters. Simulation-based calibration was prior-predictive and marginal\citep{talts2018,sailynoja2022}. The directional criteria were $P(\delta_S<-\log1.2\mid y)>0.95$, $P(\delta_\tau<-\log1.2\mid y)>0.95$ and $P(\delta_J>\log1.2\mid y)>0.95$. These correspond to ratios below approximately 0.833 or above 1.2. Any-effect was the union of individual flags. Monte Carlo uncertainty for empirical proportions used two-sided 95\% Wilson intervals, with one movie as one Bernoulli observation. No conventional null-hypothesis significance tests or multiple-testing-adjusted $P$ values were used. Coverage and decision rates were not pooled across networks evaluated on shared movies.

Timestep sensitivity used 128 independent movies at each of one, two, four, eight and sixteen internal steps per frame, separately for combined and mechanical scenarios. Ensemble means, Monte Carlo standard errors and differences relative to the finest setting were recorded. The finest setting was a numerical reference, not known exact truth. Posterior predictive checks generated 100 movies per method and selected example, conditioning on observed starting geometry and redrawing cell offsets. Featurewise interval discrepancies were descriptive and correlated. Supplementary Information gives derivations, complete numerical tables, analysis-plan provenance and the scope of each check.

\subsection*{Software, reproducibility and AI assistance}
The code used Python, NumPy, SciPy, Numba, PyTorch, ArviZ and Matplotlib; exact installed versions, random seeds, source hashes and runtime records accompany the results. Editable installation and four notebook executions were checked. The simulation and analysis scripts resume from saved records and keep exploratory and expanded datasets separate. Source data are exported as per-figure CSV files and original JSON/NPZ outputs. OpenAI Codex assisted implementation, simulation orchestration, figure construction, source checking and manuscript drafting. Figures were generated by documented numerical plotting code, and no experimental images or fabricated observations were introduced. Human verification of scientific claims, references, authorship and declarations remains pending for this draft.

\section*{Data availability}
All data in this study were generated computationally. The accompanying Source Data package contains the numerical values underlying the figures, ensemble-level records and simulation configurations. Selected trajectories and posterior draws are retained in the project outputs; remaining movies can be regenerated from stored seeds. Permanent public deposition and an accession identifier remain to be completed before submission.

\section*{Code availability}
The accompanying source snapshot contains simulator, inference, validation, plotting and manuscript-generation code, together with dependency metadata and reproduction instructions. A public repository, archival identifier and distribution license remain to be selected by the authors before submission.

@@REFERENCES@@

\section*{Acknowledgements}
To be completed by the authors; no acknowledgements have been assumed.
\section*{Funding}
Funding sources and grant information have not been supplied and require author completion.
\section*{Author contributions}
To be completed by the human authors. Author names and affiliations are intentionally blank in this draft.
\section*{Competing interests}
A competing interests declaration has not been supplied and requires author confirmation.

\clearpage
\section*{Figure legends}
\textbf{Figure 1 | A controlled framework for proximity-associated motion.}
\textbf{a}, One killer trajectory from the prespecified first strong-combined movie, colored by observed target proximity. Gray circles mark the assumed contact distance, not impermeable boundaries; the square marks the starting point. The displayed cell was chosen by maximum observed proximity variation, not recovery accuracy. \textbf{b}, Programmed contact-to-far response functions for active amplitude $S$, persistence $\tau$ and jitter $J$. The $S$ and $\tau$ curves coincide because their generating log effects are equal. \textbf{c}, The observed positions and target geometry support a likelihood-based hierarchical model and a neural summary-based posterior. \textbf{d--f}, Generating local coefficients (black), state-space posterior medians (teal) and pointwise 90\% credible bands for the same cell. Bands describe uncertainty in a shared response function evaluated along distance history. The movie contains 12 killers and 201 frames; it is one synthetic realization, not a population-average performance estimate. Source data are provided.

\textbf{Figure 2 | Calibration and information under the baseline design.}
\textbf{a}, Empirical coverage of nominal 90\% marginal intervals for six global parameters. Error bars are 95\% Wilson intervals over independent movies. Neural inference uses 1,024 new movies, state-space inference uses @@SSM_CAL_PASSED@@ accepted fits from 96 new movies, and rejection approximate Bayesian computation uses the initial 256-movie test. The dashed line marks nominal coverage. \textbf{b}, Mean posterior 90\% interval width divided by the corresponding central 90\% prior interval width; smaller values indicate more contraction. \textbf{c--e}, Generating proximity effects and posterior medians for neural (blue circles) and state-space (teal crosses) inference. The diagonal indicates exact recovery. Different ensemble sizes are retained explicitly rather than combined. Source data include movie-level results and unresolved-fit counts.

\textbf{Figure 3 | Detection of altered motion exceeds individual parameter recovery.}
\textbf{a}, Any-effect declaration rates for fixed-baseline and varied-baseline nulls. Neural denominators are 256 movies per condition; state-space denominators are @@SSM_NULL_N@@ and @@SSM_VARIED_NULL_N@@ accepted fits from 96 and 48 planned movies. \textbf{b}, Individual effect detection in the strong combined condition: 128 neural movies and @@SSM_COMBINED_N@@ accepted state-space fits from 96 planned movies. Panels \textbf{a,b} show two-sided 95\% Wilson intervals. \textbf{c}, Neural per-effect detection rates across isolated and mixed alternatives, with 128 independent movies per row. An effect is flagged when its marginal posterior probability exceeds 0.95 of crossing the stated 1.2-fold directional threshold. Any-effect is the union of these three flags. No rate is assumed to equal a nominal test size. Source data are provided.

\textbf{Figure 4 | Modeling the observation process reduces false attribution.}
The frozen fixed-noise network is compared with primary broader-noise networks that marginalize localization noise or condition on its known value. \textbf{a}, Any-effect declarations in unchanged-motility movies at localization standard deviations 0.15 and 0.60\,$\mu$m ($n=128$ per condition, shared across methods). \textbf{b}, Nominal 90\% marginal coverage over 1,024 independently simulated movies spanning noise standard deviations 0.05--0.70\,$\mu$m. \textbf{c}, Per-effect detection in 128 high-noise strong-combined movies. Error bars in \textbf{a--c} are 95\% Wilson intervals over movies. \textbf{d}, Distribution of posterior-median active-amplitude effects under the two null conditions; markers show the median across movies and bars the 10th--90th percentiles, not credible intervals or confidence intervals. Primary network seeds were fixed before evaluation. Other training seeds are shown in Supplementary Fig.~S3. These are deliberately unguarded model comparisons; the application wrapper rejects declared protocol shifts. Source data are provided.

\textbf{Figure 5 | Transfer and misspecification constrain interpretation.}
\textbf{a}, Individual effect detection for strong-combined movies generated by the supplied Euler/polarity agent-based model ($n=16$ neural; $n=4$ state-space). \textbf{b}, Any intrinsic-effect declaration under matched null, spontaneous temporal variation and mechanical-only conditions ($n=32$ per neural condition; state-space $n=32$, 8 and 8). Mechanical-only is a null for programmed intrinsic coefficients, not for observed behavior. \textbf{c}, Coverage of the zero active-amplitude proximity effect by state-space 90\% intervals for those same conditions. Bars are two-sided 95\% Wilson intervals. Low declaration rates can coexist with poor interval coverage. Small stress-test denominators and incomplete mechanical timestep convergence make these exploratory attribution checks. Source data are provided.

\clearpage
\section*{Tables}
\textbf{Table 1 | Simulation ensembles and their distinct roles.}
\small
\begin{tabular}{>{\raggedright\arraybackslash}p{.31\textwidth}>{\raggedright\arraybackslash}p{.14\textwidth}>{\raggedright\arraybackslash}p{.45\textwidth}}
\toprule
Ensemble & Movies & Purpose \\
\midrule
Original fixed-noise bank & 16,000 & Neural training and internal validation; rejection prior bank \\
Expanded noise bank & 24,000 & Matched-bank conditioned/marginalized noise training \\
New fixed-noise neural calibration & 1,024 & Independent marginal recovery; shared across four networks \\
New noise neural calibration & 1,024 & Independent recovery across noise; shared across networks \\
New state-space calibration & 96 & Independent likelihood-based recovery \\
Expanded neural nulls & $2\times256$ & Fixed and varied baseline null decisions \\
Expanded neural alternatives & $5\times128$ & Isolated, weak and strong combined changes \\
Expanded noise scenarios & $4\times128$ & Null/combined at two fixed noise values \\
Expanded state-space decisions & $96+96+48$ & Fixed null, combined and varied-baseline null \\
Original source-ABM transfer & $3\times16$ & Transfer with four state-space fits per condition \\
Original timestep study & $2\times5\times128$ & Two regimes, five internal step settings \\
\bottomrule
\end{tabular}
\normalsize

\clearpage
\section*{Figures}
\begin{center}\includegraphics[width=\textwidth]{figures/fig1_model.pdf}\end{center}
\clearpage
\begin{center}\includegraphics[width=\textwidth]{figures/fig2_calibration.pdf}\end{center}
\clearpage
\begin{center}\includegraphics[width=\textwidth]{figures/fig3_identifiability.pdf}\end{center}
\clearpage
\begin{center}\includegraphics[width=\textwidth]{figures/fig4_observation_noise.pdf}\end{center}
\clearpage
\begin{center}\includegraphics[width=\textwidth]{figures/fig5_attribution.pdf}\end{center}
\end{document}
